\section{Beyond Web5: A Substrate Map for Operator-Resident Agentic
Systems}\label{beyond-web5-a-substrate-map-for-operator-resident-agentic-systems}

\subsection{Abstract}\label{abstract}

The progression Web1 → Web2 → Web3 → Web4 is often summarized as read →
read-write → own → inhabit. Web5 has been positioned variously as
decentralized identity, as personal-server federation, or as a vague
``next layer,'' but there is little consensus on what a concrete
operator-resident substrate beyond Web4 should contain. We describe and
instrument a working substrate that orchestrates the prior four layers
into a single operator-resident shell. We identify seven primitives ---
phone-anchored decentralized identifiers (Skeet), bot-federated identity
in a biometric-gated local shell (Spacebar), signed bivector closures
for event logging (BCODE), behavioral-fidelity scoring (DECF + a
17-persona dataset, ConstellationBench), agent inheritance patterns
(OctoAgent), token-economy coordination (\$MOTION with ACODE-free /
BCODE-paid-or-earned separation), and matrix-style federation transport
(Matrix / bitchat) --- and show how they compose into a running system
with a fully tested state service (4/4 canonical ``wires'' green, 25/26
behavioral tests passing). ConstellationBench serves as an empirical
lens: we measure behavioral fidelity of frontier models across personas
and use the results to stress-test inheritance and logging primitives.
We release the dataset, runtime, and core substrate components as open
artifacts (ACODE-free), enabling other operator-resident systems to
reuse individual primitives or benchmark alternative substrate designs
against the same behavioral lens.

\section{Introduction}\label{introduction}

The last decade of web and AI infrastructure has produced a familiar
ladder: Web1 (read --- static pages, one-way information flow), Web2
(read + write --- social platforms, user-generated content), and Web3
(own --- tokenized assets, on-chain identifiers, decentralized wallets)
\citep{bernerslee1989, oconnor2023web4}. Recent proposals add two
further layers: Web4 (filter / express --- substrate-level routing,
provenance, and expression governance) and Web5 (inhabit / identify ---
operator-resident identity and control surfaces)
\citep{tbd_web5, solid_pods, ssi_survey}, but largely stop at
architectural slogans rather than concrete, measurable systems. In this
work we take a complementary approach: instead of debating the exact
definition of Web4 and Web5, we implement and instrument an
operator-resident substrate that \emph{realizes} both verbs explicitly
--- a translation-layer provenance filter (Web4) and a biometric-gated
local identity shell (Web5) --- treating Web1--3 as input surfaces. The
resulting system combines seven primitives --- identity, logging,
behavioral scoring, inheritance, incentives, and federation --- into a
working ``operator OS'' whose behavior we can probe empirically using a
17-persona benchmark (ConstellationBench) and a fully wired state
service with 4/4 health, loading, persistence, and warm-start checks
passing.

\subsection{From Web5 slogans to operator-resident
substrates}\label{from-web5-slogans-to-operator-resident-substrates}

Most Web5 discussions to date focus on protocol proposals: decentralized
identifiers (DIDs) and verifiable credentials \citep{w3c_did_core},
personal data stores and Solid pods \citep{solid_pods}, or high-level
visions of ``the web of identity'' \citep{tbd_web5}. These lines of work
are important, but they leave open a concrete systems question:
\emph{what primitives are actually required to operate a real,
operator-resident substrate that sits above existing web layers and
interacts with them in practice?}

We address this question empirically. Rather than specifying a new
protocol and leaving its deployment as future work, we start from a
running system used for day-to-day development, then factor out the
primitives that proved necessary to make it observable, auditable, and
controllable by a human operator. Concretely, our implementation exposes
seven reusable primitives:

\begin{enumerate}
\def\labelenumi{\arabic{enumi}.}
\tightlist
\item
  phone-anchored decentralized identifiers (Skeet);
\item
  a biometric-gated local control shell that federates identities
  (Spacebar v0);
\item
  signed bivector closures for event logging (BCODE);
\item
  behavioral-fidelity scoring via DECF and a 17-persona dataset
  (ConstellationBench);
\item
  agent inheritance patterns for specialized assistants (OctoAgent);
\item
  a token-economy layer for incentives (\$MOTION with ACODE/BCODE
  split);
\item
  matrix-style federation transport (Matrix / bitchat).
\end{enumerate}

Each primitive is implemented in a real codebase, exported as an open
artifact, and evaluated in the context of the running substrate.

\subsection{ConstellationBench as an empirical
lens}\label{constellationbench-as-an-empirical-lens}

ConstellationBench is a 17-persona benchmark designed to measure
behavioral fidelity of large language models along multiple DECF axes:
directionality, explicitness, coherence, and fidelity
\citep{constellationbench2026}. Rather than being the sole focus of this
paper, ConstellationBench plays the role of an \emph{empirical lens} on
the substrate: we use it to probe how different primitives behave when
populated with real personas and real models.

First, we use DECF scores across personas to quantify how consistently
frontier models reproduce intended behavioral profiles, providing a
concrete notion of ``behavioral fidelity'' that goes beyond aggregate
task accuracy. Second, we instrument the substrate so that every
evaluation run produces signed BCODE closures and serialized context,
allowing us to study how inheritance patterns (OctoAgent) and logging
interact in practice. The Phase-0 wiring --- a state service, persona
loader, session persistence, and warm-start behavior, with 4/4 canonical
wires green and 25/26 tests passing --- serves as our reproducibility
receipt: it specifies exactly which components must be operational for
other groups to replicate our experiments or to plug in alternative
primitives.

\subsection{Contributions}\label{contributions}

Our contributions are threefold:

\begin{itemize}
\tightlist
\item
  \textbf{Substrate map.} We identify seven primitives that, in
  combination, are sufficient to operate a working operator-resident
  substrate above existing web layers, and we provide an architectural
  map showing how they compose into a single local control shell.
\item
  \textbf{Empirical lens.} We use ConstellationBench as a
  behavioral-fidelity lens on this substrate, reporting results for
  multiple frontier models across 17 personas and analyzing how
  inheritance and logging primitives behave under stress.
\item
  \textbf{Open artifacts.} We release the ConstellationBench dataset,
  evaluation runtime, and core substrate components (ACODE-free) so that
  other operator-resident systems can reuse individual primitives or
  benchmark alternative designs against the same empirical lens.
\end{itemize}

\section{The Substrate Map}\label{the-substrate-map}

Rather than proposing yet another abstract ``Web5'' protocol, we start
from a running system and factor out the primitives that proved
necessary to operate it as an operator-resident substrate. Concretely,
our implementation exposes seven reusable primitives that span identity,
logging, behavioral evaluation, agent inheritance, incentives, and
transport. Table 1 summarizes these primitives, their roles in the
substrate, and the canonical implementations we release; subsequent
sections show how they compose into a single local control shell and how
their behavior can be probed empirically via ConstellationBench.

\subsubsection{Table 1: The Seven Substrate
Primitives}\label{table-1-the-seven-substrate-primitives}

{\def\LTcaptype{none} % do not increment counter
\begin{longtable}[]{@{}
  >{\raggedright\arraybackslash}p{(\linewidth - 8\tabcolsep) * \real{0.2000}}
  >{\raggedright\arraybackslash}p{(\linewidth - 8\tabcolsep) * \real{0.2000}}
  >{\raggedright\arraybackslash}p{(\linewidth - 8\tabcolsep) * \real{0.2000}}
  >{\raggedright\arraybackslash}p{(\linewidth - 8\tabcolsep) * \real{0.2000}}
  >{\raggedright\arraybackslash}p{(\linewidth - 8\tabcolsep) * \real{0.2000}}@{}}
\toprule\noalign{}
\begin{minipage}[b]{\linewidth}\raggedright
\#
\end{minipage} & \begin{minipage}[b]{\linewidth}\raggedright
Primitive
\end{minipage} & \begin{minipage}[b]{\linewidth}\raggedright
Web layer
\end{minipage} & \begin{minipage}[b]{\linewidth}\raggedright
Role in substrate
\end{minipage} & \begin{minipage}[b]{\linewidth}\raggedright
Canonical implementation
\end{minipage} \\
\midrule\noalign{}
\endhead
\bottomrule\noalign{}
\endlastfoot
1 & Phone-anchored DIDs (Skeet) & Web3 + Web5 & DID-backed identity
layer that can be bound to phone numbers, social accounts, or eSIMs;
underlies Free / Pro / Enterprise tiers & Skeet bot bindings + carrier
eSIM bindings \\
2 & Spacebar (operator-resident shell) & Web5 (inhabit / identify) &
Biometric-gated local cockpit that federates inbound identity events and
agent interactions into a single operator surface & RN+Expo native
client; humbble fork base \\
3 & BCODE (signed bivector closures) & Web4 (filter / express) &
Cryptographically signed event log; records agent actions and substrate
state transitions as bivectors with side-channel-safe serialization &
\texttt{bcode/closure.py} \\
4 & DECF + ConstellationBench & Web4 & Behavioral-fidelity evaluation:
DECF scores + 17-persona dataset provide an empirical lens for
frontier-model behavior across personas & constellation-bench-hf dataset
+ scripts \\
5 & OctoAgent (agent inheritance) & Web4 + Web5 & Encodes how
specialized agents inherit parameters, memories, and behavior from a
base operator profile while remaining auditable & OctoAgent inheritance
protocol \\
6 & \$MOTION + ACODE/BCODE split & Web3 & Incentive layer separating
free substrate code (ACODE) from paid-or-earned behavior (BCODE) and
tying it to a token economy & \$MOTION on Sonr-chain + ACODE-free
monetization model \\
7 & Matrix/bitchat federation & Web2 + Web5 & Transport mesh for
messages and events (BEAM RELAY), enabling substrate-resident agents to
interoperate across services & Matrix protocol + bitchat mesh \\
\end{longtable}
}

In §3 we use ConstellationBench (primitive 4) as an empirical lens on
the substrate, with particular attention to how primitives 3 (BCODE) and
5 (OctoAgent) behave under stress. In §4 we provide the Phase-0
reproducibility receipt covering primitives 1, 2, and 7. §5 addresses
side-channel safety and differentiation against the gatekept-enterprise
pattern (HiddenLayer) for the Web4 compute slot. §6 maps the seven
primitives to a tier ladder (Free / Pro / Enterprise) that operators
traverse as they grow from social-account auth to carrier-attested
identity. §7 discusses PACT --- a protocol-layer contribution emerging
from the substrate's identity-aware-authorization needs --- and future
verticals.

\section{Empirical motivation: what ConstellationBench actually
shows}\label{empirical-motivation-what-constellationbench-actually-shows}

The substrate map in §2 is architectural; this section explains why
those particular pieces and abstractions are necessary at all.
ConstellationBench provides the empirical lens: instead of scoring
models only on generic task success, it measures two orthogonal
dimensions that matter for an operator-resident system --- did the model
solve the problem, and did it behave like the specific persona it
claimed to be. Using that lens across a variety of models and
orchestration patterns yields three consistent, somewhat uncomfortable
findings: (i) ``budget'' models can beat frontier models on persona
fidelity, (ii) a non-speaking audience agent can materially improve
quality and fidelity at zero token cost, and (iii) multi-agent teams
beyond a small size tend to collapse toward generic behavior rather than
diversifying. These results are what drive the substrate toward explicit
persona metrics, audience-aware orchestration, and a protocol layer that
treats identity and trace as first-class objects rather than comments in
a prompt.

\subsection{Setup and evaluation
lens}\label{setup-and-evaluation-lens}

Experiments use ConstellationBench, a persona-sensitive benchmark that
scores model outputs along two orthogonal axes: task quality (did the
answer solve the problem?) and persona fidelity (did it sound and behave
like the specified agent, not like a generic assistant). Models are
evaluated under the same prompt scaffolds, tools, and scoring rubric to
keep comparisons reproducible rather than anecdotal. Each condition is
run across a panel of personalities so that results reflect structural
effects, not idiosyncrasies of one archetype
\citep{constellationbench2026}.

\subsection{Headline result: budget models beat frontier on
persona
fidelity}\label{headline-result-budget-models-beat-frontier-on-persona-fidelity}

The first result is counter-intuitive if one only tracks leaderboard
scores: smaller ``budget'' models (Anthropic Haiku and Sonnet)
outperform frontier-class models (Opus, GPT-4-class) on persona fidelity
by roughly 8 percentage points while achieving comparable task quality.
In other words, cheaper models stay ``in character'' more reliably, even
though larger models score higher on generic benchmarks.

Two features make this result publication-grade rather than a one-off
anomaly:

\begin{itemize}
\tightlist
\item
  The effect is consistent across multiple personas and tasks, not
  driven by a single outlier character.
\item
  The experiment is inexpensive to reproduce: the headline
  budget-vs-frontier subset costs on the order of \$60 of API spend
  (§8.2 details a \$115 full protocol covering all 22 models × 17
  personas × 7 benchmarks); other groups can verify or refute either
  tier without bespoke infrastructure.
\end{itemize}

The practical implication is that fidelity is not a simple monotone
function of model scale. Once a model has been heavily trained and
reinforced to act like a safe, polite, generic assistant, adding more
parameters and training steps appears to strengthen that generic voice
faster than it strengthens the capacity to inhabit distinct personas.
The substrate therefore cannot assume ``bigger is better'' for agentic
work; it must be able to route to cheaper models when fidelity, not raw
capability, is the primary objective.

\subsection{The alignment trap: generic safety vs specific
identity}\label{the-alignment-trap-generic-safety-vs-specific-identity}

The fidelity gap points to a broader pattern the experiments keep
surfacing: what can be called an ``alignment trap.'' Modern frontier
models are aligned, in the safety sense, to be inoffensive, broadly
helpful, and stylistically uniform. That is appropriate for
single-assistant experiences, but it comes at a cost for multi-agent or
persona-specific systems.

Empirically, as alignment pressure and generic safety training increase,
outputs converge toward a narrow band of ``default assistant'' behavior
even when the prompt explicitly specifies a different character. In
ConstellationBench runs, this shows up as:

\begin{itemize}
\tightlist
\item
  High task scores but middling persona scores for frontier models under
  strong safety scaffolds.
\item
  Higher persona scores for smaller models or less-aligned checkpoints
  under identical persona instructions.
\end{itemize}

The alignment trap, then, is the regime in which marginal safety
training continues to improve generic assistant behavior but actively
erodes the system's ability to realize distinct, inspectable personas. A
substrate that cares about behavioral fidelity must treat alignment as a
multi-objective problem: keep safety constraints, but do not allow them
to collapse the personality space to a single point.

\subsection{Audience effect: the ``silent Guardian'' as
stabilizer}\label{audience-effect-the-silent-guardian-as-stabilizer}

A second robust effect is the ``audience'' or ``silent Guardian''
phenomenon. Adding an additional, non-generating agent to the system
prompt --- a named reviewer or overseer that never speaks to the user
but is described as evaluating the assistant's work --- produces a
measurable improvement in output quality and persona fidelity, despite
consuming no extra tokens at inference time.

Across runs, the presence of a silent Guardian agent yields
approximately:

\begin{itemize}
\tightlist
\item
  +14\% relative lift in task quality.
\item
  +10.7\% relative lift in persona fidelity.
\end{itemize}

Mechanistically, the effect looks less like classical self-critique and
more like social priming. The base model is prompted to treat its own
output as something that will be watched and graded by a peer or
supervisor, and its behavior shifts accordingly. The key is that the
Guardian never generates text or calls tools; it exists only as a
parametric audience in the system prompt.

For the substrate, this suggests that some of the benefits usually
attributed to more complex multi-agent architectures can be achieved by
simpler orchestration patterns: a single generating agent plus a
structured ``social context'' in the prompt. It also hints that
evaluation is not only something the substrate does after the fact; the
mere presence of an internal audience can change the distribution of
outputs at generation time.

\subsection{Team size and ``reef bleaching'': when more agents
hurt}\label{team-size-and-reef-bleaching-when-more-agents-hurt}

A third family of experiments probes how performance changes as the
number of cooperating agents increases. Intuitively, one might expect
larger teams to be strictly better: more specialist roles, more tools,
more perspectives. In practice, the runs show a saturation point and
then a decline.

As agent count increases past a small team --- roughly three to four
specialized roles --- the ensemble begins to exhibit:

\begin{itemize}
\tightlist
\item
  Longer, more redundant traces with little incremental improvement in
  task quality.
\item
  Blurring of persona boundaries: agents start to sound more alike, and
  their supposed specializations collapse toward a shared, generic
  style.
\item
  Increased chances of contradiction or oscillation as agents rewrite
  each other's work without a strong coordinator.
\end{itemize}

This pattern is described as ``reef bleaching'': the coral (distinct
personas and viewpoints) lose color as the environment becomes too hot
and crowded, leaving behind a pale, homogeneous structure. The
implication is that agent orchestration has a non-trivial optimum;
beyond a certain point, adding more agents weakens the very diversity
and specialization they were introduced to provide.

For the substrate, this argues in favor of small, well-defined teams
with clear coordinator roles, rather than arbitrarily large swarms. It
also reinforces the need for behavioral metrics that can detect
bleaching --- e.g., measuring how distinguishable agents remain over
time --- rather than relying solely on aggregate task scores.

\subsection{Protocol implications: a forward reference to
PACT}\label{protocol-implications-a-forward-reference-to-pact}

Taken together, these findings suggest that the protocol layer itself
must carry more than tool calls and raw messages: it must represent who
is acting, how they are allowed to act, under what audience, and with
what trace. Section 7 formalizes the PACT envelope as a typed operator
on session state. Here we only note that every benchmark call in
§§3.2--3.5 already flows through PACT → BCODE → sidecar, so the
empirical results above are not contingent on a new protocol --- they
are produced by one.

\section{Phase-0 system execution receipt: bringing up the state
service}\label{phase-0-system-execution-receipt-bringing-up-the-state-service}

The substrate argument only matters if independent readers can stand up
the same operational baseline. Before introducing credits, PACT, or
multi-agent orchestration, the stack therefore went through a
constrained ``Phase 0'' focussed entirely on bringing up a minimal,
inspectable state service and proving that it behaves deterministically
under a fixed set of tests. That Phase-0 snapshot functions as the
execution receipt for all higher-level claims --- distinct from the
benchmark replication recipe in §8.2, which addresses a different
question (rerunning the ConstellationBench dataset).

\subsection{Scope and design of Phase
0}\label{scope-and-design-of-phase-0}

Phase 0 is intentionally narrow. It does not attempt to implement
credits, Skeet identity, or tiering. Instead, it defines four canonical
``wires'' that any operator-resident substrate must satisfy:

\begin{itemize}
\tightlist
\item
  \textbf{V1 State Service Health:} the service responds on a known port
  with a health check.
\item
  \textbf{V2 Profile Loading:} personas can be loaded from disk into
  memory via an HTTP interface.
\item
  \textbf{V3 Session Persistence:} sessions survive across calls and
  restarts in a controlled way.
\item
  \textbf{V4 Warm Start:} pre-seeded sessions can be resumed with their
  prior context.
\end{itemize}

The stack for this service is deliberately conservative: FastAPI +
SQLAlchemy 2.0 + Postgres 16 on the backend, with configuration
centralized in an \texttt{airlock-coordination/config.py} module that
sets \texttt{STATE\_SERVICE\_HOST}, \texttt{STATE\_SERVICE\_PORT},
\texttt{STATE\_SERVICE\_BASE}, and a fixed
\texttt{ONBOARDING\_MAX\_TURNS=10} constant used by downstream agents.
No exotic dependencies are introduced; everything that runs in Phase 0
can be deployed on a single developer laptop.

\subsection{Implementation details and wire
tests}\label{implementation-details-and-wire-tests}

Implementation proceeds in three commits that are explicitly documented
in the runbook:

\begin{itemize}
\tightlist
\item
  \textbf{Phase 0.4 (verification only):} confirm that a ``19th plan''
  file remains in the tree because it is referenced in
  \texttt{AGENTS.md:22}; no code changes, but establishes that prior
  planning material is preserved rather than silently deleted.
\item
  \textbf{Phase 0.5 (\texttt{1004b4c}):} reset \texttt{.gitignore} to a
  known good baseline, extending it to include \texttt{reports/*.yaml},
  \texttt{relays/repos/*.json}, \texttt{relays/personas/*.json}, and
  placing \texttt{.gitkeep} markers where appropriate; add a README note
  describing a \texttt{/refresh} endpoint that will later be used by
  higher-level components.
\item
  \textbf{Phase 0.6 (\texttt{415df10}):} create \texttt{config.py} with
  the state-service host/port/base settings and
  \texttt{ONBOARDING\_MAX\_TURNS=10}, patch \texttt{server.py} and
  \texttt{validate-wires.py} to import configuration rather than
  hard-coding values; synchronize these ports with the runbook.
\item
  \textbf{Phase 0.7 (runtime):} create a dedicated virtual environment,
  install \texttt{fastapi}, \texttt{uvicorn}, \texttt{pydantic},
  \texttt{python-ulid}, and \texttt{pyyaml}, freeze the resulting
  \texttt{requirements.txt}, start the service on port 8100, and run
  \texttt{validate-wires.py}. At this snapshot the service passes 24 out
  of 26 checks; all four canonical wires are green, with one
  non-blocking V5 sub-check (\texttt{user\_profile} field on an
  auxiliary endpoint) flagged for later work.
\end{itemize}

The wire harness is intentionally simple. It issues HTTP requests to
\texttt{/health} (V1), \texttt{/persona/profiles} (V2), and
\texttt{/persona/warm-start} (V4), and it exercises a short
create-read-update cycle on session endpoints to verify V3 persistence.
Success is defined as all four wires returning expected shapes and
status codes; any deviation fails the test suite.

\subsection{Failure modes and
forensics}\label{failure-modes-and-forensics}

Phase 0 also documents the first real failure the stack encountered: a
disk-full event that truncated \texttt{state.json} and
\texttt{queue.json}. Instead of silently repairing and moving on, the
runbook records the corruption, notes that both files were reset to
valid empty shapes, and preserves the originals as
\texttt{state.json.corrupt-20260429} for inspection. This is not a
research contribution in itself, but it demonstrates the intended
posture: the substrate is expected to keep forensic receipts of its own
failures, not only its successes.

Similarly, Phase 0 discovers two structural flags that are explicitly
carried forward rather than patched over:

\begin{itemize}
\tightlist
\item
  \texttt{constellation.py} currently expects \texttt{reports/*.json} on
  disk, but the repository contains 62 \texttt{reports/*.yaml} files
  from a prior run. The mismatch is acknowledged and deferred to
  post-Phase-0 investigation.
\item
  \texttt{airlock-persona} is symlinked from a nested
  \texttt{forge-app/Forge\ UI/airlock-persona} path; the runbook notes
  that this should eventually move behind an \texttt{OTTO\_PERSONA\_DIR}
  environment variable for portability.
\end{itemize}

In both cases the Phase-0 snapshot includes the inconsistency and a
pointer to its resolution path, so that readers attempting to reproduce
the environment can see the system as it actually was rather than as a
cleaned-up ideal.

\subsection{Why this matters for the substrate
argument}\label{why-this-matters-for-the-substrate-argument}

From the standpoint of this paper, Phase 0 plays a specific role. It
provides:

\begin{itemize}
\tightlist
\item
  A pinned \textbf{environment}: a known port (\texttt{8100}), a fixed
  configuration file, a frozen \texttt{requirements.txt}, and a small
  number of commits that define the state of the repository when the
  tests passed.
\item
  A minimal but sufficient \textbf{API surface}: health, persona
  loading, warm-start, and session persistence, all exercised by a
  public \texttt{validate-wires.py} script.
\item
  A concrete \textbf{failure history}: corruption and format mismatches
  that future work must respect, not erase.
\end{itemize}

All later sections --- credits, Skeet identity, PACT, and the
Web3/Web4/Web5 ladder --- assume that some state service exists to hold
persona definitions and session context. Phase 0 is the evidence that
such a service can be implemented in a way that is both straightforward
to run and precise enough to serve as a baseline. A reader who wishes to
reproduce the substrate's behavior can start here: check out the
repository at the documented commits, create the venv, run the service
on port 8100, and confirm that the four wires are green before
attempting to replicate ConstellationBench results or introduce more
complex orchestration.

\section{Safety, side-channels, and the Web4 security
layer}\label{safety-side-channels-and-the-web4-security-layer}

Modern AI systems are increasingly deployed behind dashboards that
advertise ``glass-box'' observability. In practice, most of these
surfaces reason about scalar metrics (latency, cost, error rate) while
leaving a large class of structural leaks and behavioral commitments
unmodeled. This section argues that a substrate intended to sit
underneath Web3--Web5 needs two things the current stack largely lacks:
(i) explicit side-channel discipline around numeric telemetry, and (ii)
a security layer that treats behavioral commitments the way cryptography
treats hash functions and signatures.

\subsection{Numeric telemetry as a
side-channel}\label{numeric-telemetry-as-a-side-channel}

A recent public incident around model-usage telemetry illustrates how
easily an apparently benign field can become a side-channel. In that
case, a provider's billing endpoint streamed unrounded IEEE-754 doubles
such as \(0.16327272727272726\) for ``used/limit'' credit ratios
\citep{katsarosgiannis2026}. Because both numerator and denominator live
in bounded integer ranges, an attacker can apply the Stern--Brocot tree
to recover the simplest fraction within that float's rounding bucket,
yielding exact \((\text{used},\text{limit})\) pairs. Aggregating
multiple observations and taking least common multiples of recovered
denominators reveals subscription caps and tier boundaries that were
never explicitly documented.

The same mathematics applies to any substrate that emits full-precision
floats derived from bounded integers. BCODE ``truth signals'' are
defined as per-call cost estimates for a given closure: predicted
compute cost, verification overhead, and ledger append cost are combined
into a scalar that can easily land at values like 0.0037 USD. If such
values are exposed verbatim over an API --- whether for transparency or
internal tooling --- they can leak exact price curves, internal
efficiency margins, and cross-surface allocation rules in precisely the
way the usage example does.

The substrate therefore adopts a simple but strict rule: all externally
visible floats tied to bounded integer quantities are rounded at the
serialization boundary. For BCODE truth signals, the internal cost model
may compute a true cost of 0.0037, but the operator is charged an
integer multiple of a substrate unit (for example, 0.01 USD), and only
that rounded figure is exposed. The difference between true and rounded
cost is accumulated in an operator-scoped ``rounding surplus'' treasury
that funds a self-contained promotion (e.g., ``buy four calls, the fifth
is effectively paid by surplus'') without requiring any cross-subsidy
from unrelated revenue. For quotas and usage, the substrate returns
separate integer \texttt{used} and \texttt{limit} fields or a percentage
rounded to two decimal places, never the raw ratio; DECF scores and
bivector components are rounded to three and four decimals respectively
before being logged.

This is not just an accounting convenience. It plays three roles that
matter for safety:

\begin{itemize}
\tightlist
\item
  It \textbf{caps the information capacity} of telemetry fields so that
  simple rational reconstruction attacks cannot recover internal
  parameters.
\item
  It keeps the \textbf{cooperative audit ledger integer-clean}: every
  BCODE closure row carries whole-unit charges and whole-unit surplus,
  eliminating dust that complicates replay and reconciliation.
\item
  It creates a tunable \textbf{``humid dial''}: by choosing how
  aggressively to round (to the nearest cent, the nearest 0.10, or
  0.25), the operator can trade off surplus accumulation and promotional
  capacity against operational variance, with each dial turn recorded as
  a signed configuration change in the truth manifest.
\end{itemize}

In other words, the same rounding step that closes a side-channel also
finances operator-visible promotions and simplifies ledger replay. This
sort of multi-role primitive --- security, economics, and
reproducibility at once --- is characteristic of the substrate's design.

\subsection{Behavioral commitments and layered
trust}\label{behavioral-commitments-and-layered-trust}

Side-channels are only one part of the safety story. A substrate that
aspires to carry behavioral commitments --- the promise that ``this
session will not be used to reconstruct you later'' --- needs a design
discipline closer to modern cryptography than to typical SaaS telemetry.
A companion research note on ``SHA-256 as a layered trust object''
provides the template \citep{wp3_sha256_layered_trust}.

That note decomposes any serious trust claim into five layers:

\begin{enumerate}
\def\labelenumi{\arabic{enumi}.}
\tightlist
\item
  \textbf{Primitive} --- the mathematical object itself (e.g., a hash
  function, or here a behavioral commitment function).
\item
  \textbf{Attack catalog} --- known structural and probabilistic attacks
  against the primitive (reduced-round breaks, partial preimages,
  attribute inference).
\item
  \textbf{Protocol context} --- how the primitive is used (for
  passwords, for proof-of-work, or in this case for session anchoring
  and entanglement-safe handshakes).
\item
  \textbf{Implementation} --- concrete code and hardware, including
  side-channel properties.
\item
  \textbf{Operational context} --- who runs it, under what threat model,
  with what audits and regulatory constraints.
\end{enumerate}

Applied to a behavioral substrate, this frame yields a clear research
agenda. At Layer 1, a candidate behavioral commitment function must be
specified --- a mapping from a user's initial behavioral state (for
example, a DECF vector and optional embeddings) to a non-invertible
commitment that the vendor can verify but not reconstruct. At Layer 2,
an attack catalog must answer questions analogous to cryptographic
cryptanalysis: given a commitment, can an adversary reconstruct the
underlying drives (preimage), infer specific attributes (attribute
inference), or construct three-way relationships that violate claimed
independence (3XOR-style attacks)? At Layers 3--5, the non-separability
and ``entanglement-safe handshake'' papers define protocol topologies in
which such a primitive would be deployed
\citep{wp2_entanglement_safe_handshake}, while the present stack begins
to specify implementation and operational behavior via PACT envelopes
and BCODE audits \citep{pact_protocol_spec}.

The crucial point is that a ``glass-box'' claim is incomplete if it
addresses only Layers 4 and 5 (implementation and ops) without
specifying a primitive and its attack surface. By explicitly tying PACT
and BCODE to a layered trust frame, the substrate creates space for a
future behavioral commitment primitive to be dropped in and analyzed
with the same seriousness SHA-256 has received over two decades.

\subsection{HiddenLayer and the Web4 security
slot}\label{hiddenlayer-and-the-web4-security-slot}

Commercial products already occupy a ``hidden-layer security'' slot for
AI deployments. Vendors like HiddenLayer \citep{hiddenlayer} position
themselves as SOC-integrated monitors that scan models for adversarial
attacks, exfiltration attempts, and anomalous behavior, often with
integrations into existing logging and SIEM systems. In the taxonomy
used here, this is a Web4-tier offering: a compute-adjacent layer that
observes but does not own the orchestration surface.

The substrate takes a different stance. Rather than bolting a
third-party detector onto an otherwise opaque stack, it treats security
and observability as first-class citizens within the compute layer
itself. Every scored interaction in ConstellationBench produces a BCODE
closure (defined formally in §7.1) that binds persona drives, model
identity, fidelity signal, and historical context into a single object.
Each closure is serialized as JSON, signed (with a stub key in the
benchmark release and an ed25519/Sonr DID key in production), and stored
in a sidecar logfile where replay verifiers can recompute and verify it.

Practically, this means that tasks a HiddenLayer-style service would
perform --- detecting off-policy behavior, spotting unusual error
payloads, ensuring that audit trails cannot be tampered with --- are
instead implemented as substrate features:

\begin{itemize}
\tightlist
\item
  \textbf{Untangler gates} refuse to accept closures whose magnitude or
  direction deviates too far from a configured norm, blocking ``weird''
  interactions even if the upstream model produced a syntactically valid
  response.
\item
  \textbf{Layered trust} is embedded in the protocol: each PACT call
  carries a hash of its payload, a pointer to its closure, and zone
  metadata (green / yellow / red) that can be independently audited at
  the SHA-256 structural-literature level rather than treated as
  application logs.
\item
  \textbf{Side-channels are addressed at the primitive level}: rounding
  rules, treasury mechanics, and commitment designs are documented and
  testable, rather than left as implementation details behind a billing
  API.
\end{itemize}

HiddenLayer's existence is thus treated as empirical validation that the
Web4 security slot is economically real, not as a competitor to be
displaced. The substrate's contribution is to show that the same class
of guarantees --- attack visibility, anomaly detection, tamper-evident
logs --- can be expressed as open, signed, replayable artifacts (PACT
envelopes and BCODE closures) that live inside the system rather than at
its perimeter.

Taken together, these design choices support the broader thesis of the
paper: that if behavior is fundamentally non-separable, the only honest
way to deploy AI at scale is to let operators see --- and, when needed,
fork --- the layer that watches their models, not just the layer that
serves their responses.

\subsection{Provenance as substrate governance: the Web4
filter}\label{provenance-as-substrate-governance-the-web4-filter}

The post-hoc detection problem (``is this text AI-generated?'') is
unsolvable at internet scale because aggressive filters destroy human
data and loose filters admit synthetic contamination. Our substrate
avoids that trap by relocating the filter from the web-scrape layer to
the translation layer between ACODE (open, application-facing) and BCODE
(signed, audited substrate-native), where the router already knows
session context, identity, and operation type. Each PACT envelope's
audit-trace binding includes a structured provenance block ---
\texttt{origin\_type}, \texttt{origin\_actor},
\texttt{translation\_mode}, \texttt{trust\_class},
\texttt{retain\_policy} --- tagged at the moment of passage. This
converts the impossible global inference problem into a tractable local
governance problem: the substrate does not detect AI content on the
internet; it tags what kind of entity is producing this payload as it
crosses into the audited layer.

A small policy engine consumes the provenance block and returns one of
three router decisions --- \texttt{allow},
\texttt{allow\_with\_downgrade}, or \texttt{deny} --- based on six
rules. The sovereign-human path (\texttt{origin\_type=human},
\texttt{trust\_class=sovereign}) is the only class eligible for default
elevation to \texttt{full\_audit} retention; agent-originated material
defaults to \texttt{closure\_only} retention even when the agent acts on
a sovereign user's behalf, preventing silent trust upgrades.
Synthetic-only inputs cannot enter \texttt{full\_audit} raw retention
but may participate in derived operations and bounded scoring. Mixed
inputs are downgraded to \texttt{closure\_only} in high-integrity zones.
Missing provenance is a protocol error: requests cannot cross ACODE into
BCODE without it. Chambers and zones can impose stricter overrides ---
for example, a \texttt{Control} chamber may forbid synthetic origins
from triggering mutating operations entirely.

The architectural move is what matters: this is the canonical Web4
mechanism in our substrate. Where commercial products like HiddenLayer
(§5.3) operate at the model-monitoring perimeter, our provenance filter
operates at the message-passing interior --- at the precise moment a
payload becomes a candidate for durable signed storage. The filter is
not a model; it is a typed policy applied to a typed envelope. This
makes it amenable to formal verification, replay-auditable across
substrate instances, and (because it is ACODE-free) inspectable and
forkable by operators who want stricter or alternative governance rules.
The full implementation is released as a Pydantic reference module so
that Otto-class router work --- including the in-process ExecutionRouter
formalized in §7.1 --- can consume it directly.

\section{Tier ladder: coupling identity assurance and compute
governance}\label{tier-ladder-coupling-identity-assurance-and-compute-governance}

A substrate that claims to be ``operator-resident'' must expose its
resource constraints in a way that individual operators can understand
and audit. Rather than treating identity verification and compute limits
as separate concerns, the current stack couples them into a single tier
ladder. Each rung corresponds to (i) a specific identity surface with a
clear assurance level, and (ii) a compute envelope expressed in
off-chain credits with on-chain \$MOTION reserved for incentives.

\subsection{Constellation Credits as the compute
primitive}\label{constellation-credits-as-the-compute-primitive}

On the compute side, the substrate uses an internal accounting unit ---
Constellation Credits (CC) --- to meter model calls, storage, and BCODE
closure operations. CCs are held in a Postgres ledger keyed by operator
DID, with each entry representing a discrete event (e.g., ``one
Sonnet-grade call with corkscrew orchestration'' or ``one Haiku-grade
judge pass''). The credit design separates three concerns:

\begin{itemize}
\tightlist
\item
  \textbf{Unit pricing:} each action has a fixed CC cost derived from
  its expected USD price at a reference provider (OpenRouter) and a
  small safety margin.
\item
  \textbf{Budget enforcement:} per-operator and global pool caps are
  enforced via simple invariants (no negative balances, nightly pool
  limits) rather than complex dynamic pricing.
\item
  \textbf{Auditability:} every CC debit or credit emits a BCODE closure
  with economic fields (charged amount, rounding surplus, treasury
  credit) so that cost flows can be replayed and verified.
\end{itemize}

Crucially, CCs are \textbf{off-chain}. They are not a token; they do not
live on a blockchain; they are not intended to be traded. The on-chain
asset, \$MOTION, is reserved for ecosystem incentives and external
interoperability. A later bridge (Phase 4 in the pricing documents) will
allow operators to convert \$MOTION into CC and vice versa at defined
rates, but the two remain distinct: CC is ``gas'' inside the substrate,
\$MOTION is a public asset outside it.

This separation allows the substrate to offer predictable,
dollar-denominated pricing for operators while still supporting on-chain
mechanics where they make sense (for example, rewarding OSS
contributions with \$MOTION that can be converted into CC). The full
Constellation Credits design documents are held internally during
economics finalization (preserving the ACODE/BCODE split while the
\$MOTION bridge is finalized); the public-facing implementation will
follow once the on-chain mechanics and any required regulatory clearance
are settled, and the artifact map in §8.1 reflects this status.

\subsection{Skeet identity tiers as the identity
primitive}\label{skeet-identity-tiers-as-the-identity-primitive}

On the identity side, Skeet defines three increasingly strong identity
tiers, each corresponding to a different ``auth surface.''

\begin{itemize}
\tightlist
\item
  \textbf{Free (ACODE-free):} operators bind an existing social account
  (Discord, Telegram, later WhatsApp) via OAuth. A bot in their direct
  messages becomes the auth surface, receiving login codes and approval
  prompts. No phone number is required; the substrate verifies control
  of two accounts (the device making the request and the social account
  receiving the DM) as a built-in second factor.
\item
  \textbf{Pro (BCODE-paid or earned):} operators bind a real carrier
  eSIM, using services such as Noble Mobile or Mint. The substrate
  verifies via a number-lookup service that the line is carrier-issued
  (not VoIP) and binds that hardware-backed identity to the operator's
  DID. This tier is priced around USD 19.99/month or equivalent
  \$MOTION, and is intended for privacy-serious users and flows where
  services reject VoIP or social-only auth.
\item
  \textbf{Enterprise / PSA-as-carrier (future):} the substrate itself
  issues eSIMs, enabling family meshes, strong Sybil resistance, and
  higher-assurance flows. This tier sits outside the current
  implementation but is documented as a long-term direction.
\end{itemize}

From an architectural perspective, the Free tier is ACODE-level: the bot
client and OAuth templates are open and freely forkable, with no
per-message cost beyond the platforms' own hosting. Pro and Enterprise
tiers are BCODE-level: they integrate paid carrier infrastructure,
require non-trivial operational work, and thus belong behind a pricing
wall.

\subsection{Bundling identity and compute into a single
ladder}\label{bundling-identity-and-compute-into-a-single-ladder}

Rather than present operators with two separate ladders (``how verified
are you?'' and ``how many credits do you have?''), the substrate bundles
them. The Constellation Credits design documents sketch a mapping of the
form:

\begin{itemize}
\tightlist
\item
  \textbf{Tier 0 --- Free:} social-account Skeet identity; a small
  monthly CC allocation suitable for light experimentation (for example,
  tens of ConstellationBench-scale calls), funded by ACODE-free
  mechanisms and introductory promotions.
\item
  \textbf{Tier 1 --- Pro:} carrier-backed Skeet identity; a larger
  monthly CC allocation (on the order of thousands of credits) plus a
  quota of high-end model invocations; priced in USD or \$MOTION;
  suitable for serious individual use.
\item
  \textbf{Tier 2 --- BYOK Pro:} operators supply their own OpenRouter or
  vendor API keys; CC is used only for BCODE and PACT overhead, with
  model costs billed directly; identity can be Free or Pro; intended for
  teams and integrators.
\item
  \textbf{Tier 3 --- Enterprise:} organizational Skeet identity and eSIM
  provisioning; custom CC envelopes and governance; intended for
  institutions with their own risk and compliance regimes.
\end{itemize}

Two properties of this design are worth underscoring.

First, \textbf{compute is governed at the protocol level}, not just at
the UI. PACT envelopes carry not only persona and zone information but
also per-call cost ceilings and remaining CC balances. A green-zone
Haiku evaluation call cannot silently trigger a red-zone Opus
re-amplification that would blow through a nightly budget; the protocol
rejects such escalations unless an explicit approval envelope raises the
ceiling.

Second, \textbf{pricing logic is expressed as BCODE}, not as opaque
billing code. The ``truth signal rounding'' mechanism defines how
per-call true costs are rounded up to a substrate unit (e.g., 0.01 USD),
how the surplus accumulates in an operator-bound treasury, and how every
fifth call is funded from that treasury to create a ``buy four, get one
free'' promotion without cutting margin. Each of these steps produces a
signed BivectorRecord with economic fields that can be replayed; dial
changes (for example, rounding to 0.10 instead of 0.01 for more
aggressive promotions) are stored in a versioned ``truth signal dial''
table governed by the same approval gates as other BCODE changes.

The result is a \textbf{single ladder} that operators can reason about:

\begin{itemize}
\tightlist
\item
  At the bottom, anyone with a social account and an email can obtain a
  DID, a bot-mediated auth surface, and a small CC allocation for free.
\item
  As they climb, they add stronger identity (carrier eSIM), larger and
  more flexible compute envelopes, and additional features (e.g., Skeet
  keyrings, Chainlink-backed cross-chain key attestations) without
  losing the ability to audit what the substrate is doing on their
  behalf.
\end{itemize}

From a research perspective, the important contribution is not the
specific price points, which will drift, but the coupling: identity
assurance and compute governance are treated as two facets of the same
object, and that object is represented in the same substrate-native way
as behavioral events --- via PACT envelopes and BCODE closures that can
be inspected, replayed, and, if necessary, forked.

\section{Discussion: PACT, non-separability, and protocol-level
guarantees}\label{discussion-pact-non-separability-and-protocol-level-guarantees}

The preceding sections treated the substrate as an engineering artifact:
a benchmark, a set of services, a tier ladder. This section steps back
and describes it as a dynamical system. The goal is not to propose a new
formalism for its own sake, but to show that two choices --- using PACT
envelopes as the unit of interaction and treating behavior plus context
as a non-separable state --- are structurally necessary if the system is
going to support the interference and phase-transition phenomena
observed in prior work.

\subsection{PACT as a typed operator on session
state}\label{pact-as-a-typed-operator-on-session-state}

PACT --- the \textbf{Protocol for Agents, Context, and Trace} --- is the
substrate's protocol-layer object. Concretely, every interaction with
the substrate takes the form of a PACT envelope: a structured record
that binds, in the same object, (i) identity (operator DID, signing
key), (ii) persona and role context (DECF profile, zone, tier, current
chamber/view constraints), (iii) autonomy gates (which tools and
escalations are permitted under this envelope), (iv) budget ceilings
(per-call cost limit, remaining credits), (v) zone metadata
(green/yellow/red risk classification), and (vi) a durable audit trace
(links to prior closures, signed transition history). It is tempting to
treat this as ``just JSON.'' For the purpose of analysis, it is more
useful to view PACT as a typed operator on a product space.

Let \(X\) denote the space of request headers, including:

\begin{itemize}
\tightlist
\item
  persona id, zone, tier, and routing hints
\item
  per-call budget ceilings and remaining credit
\item
  session identifiers and turn counters.
\end{itemize}

Let \(Y\) denote the space of model-level outcomes (response text, log
probabilities if available, timing, and error flags), and \(Z\) the
space of closure state as defined in the BCODE operator:

\[
C = P \wedge L \oplus R \oplus W,
\]

where \(P\) is the persona 4-vector, \(L\) the SHA-256-derived model
4-vector, \(R\) the response-distributed fidelity vector, and \(W\) a
rolling history vector. A single PACT application can then be written as
an operator

\[
\Pi : X \to X \times Y \times Z,
\]

which maps an input header \(x \in X\) to a triple \((x', y, z)\), where
\(x'\) is the updated header, \(y\) the model outcome, and \(z = C\) the
resulting closure.

The protocol defines an acceptance band over closures using the
Frobenius norm on the grade-2 part of \(C\). Let \(\|\cdot\|_F\) denote
this norm and \(d > 0\) a configured threshold. A closure is accepted if
\(\|C\|_F \le d\); otherwise it is emitted with
\texttt{closure\_status\ =\ rejected} but still written to the audit
log. This yields a simple but important property.

\begin{quote}
\textbf{Proposition 7.1 (Closure-band stability).} For any fixed persona
\(P\) and model family in the benchmark protocol, there exists a
threshold \(d > 0\) such that all accepted closures satisfy
\(\|C\|_F \le d\), and the replay verifier will detect any post-hoc
perturbation that pushes \(\|C\|_F\) outside this band, even if the JSON
remains syntactically well-formed.
\end{quote}

Operationally, this is implemented by the \texttt{verify.py} module: for
each line in the signed sidecar file, the verifier recomputes \(C\) from
stored \(P, L, R, W\), checks the norm against the acceptance band, and
reports mismatches separately from signature failures. Conceptually, it
elevates closures from ``logs'' to typed objects with a well-defined
acceptance predicate. PACT is not merely a request/response wrapper; it
is an operator that evolves session state in a space where certain
invariants --- closure norm bounds, budget constraints, and zone
restrictions --- are enforced.

Budget is treated similarly. Let \(B_t\) denote the remaining
Constellation Credit balance at step \(t\), as recorded in the header
component of the PACT envelope. For any sequence of PACT applications
\(\Pi_t\), the protocol enforces

\[
B_{t+1} \le B_t
\]

with equality only for purely observational steps that do not trigger
model compute. Together with the closure norm band, this defines a pair
of monotone quantities: credit is non-increasing, and the cumulative
audit trail (the multiset of closures) is non-decreasing. PACT thus
defines a discrete-time dynamical system on \((X, B, C)\) with explicit
guards on both economics and behavior.

\subsection{Non-separability and the bivector
interlingua}\label{non-separability-and-the-bivector-interlingua}

The PACT operator describes how state evolves. The second question is
what form that state should take. Prior work on topological intelligence
argued that behavioral identity is not a property of individual models
but of a network, and that quantum-cognition-style interference effects
appear when multiple persona contexts are composed. Those findings
suggest that any substrate which claims to reason about behavior must
treat behavior and context as a non-separable object, not as independent
axes routed in isolation.

Formally, let \(\mathcal{H}_{\text{behavior}}\) be a real vector space
spanned by DECF basis vectors and their derived components, and let
\(\mathcal{H}_{\text{context}}\) be a vector space spanned by context
features: model family, tier, zone, historical window, and external
signals such as observation framing. A combined session state lives in
the tensor product

\[
\Psi \in \mathcal{H}_{\text{behavior}} \otimes \mathcal{H}_{\text{context}}.
\]

A state is \textbf{separable} if there exist
\(b \in \mathcal{H}_{\text{behavior}}\) and
\(c \in \mathcal{H}_{\text{context}}\) such that \(\Psi = b \otimes c\).
It is \textbf{non-separable} if no such factorization exists and
\(\Psi\) can only be written as a sum

\[
\Psi = \sum_{i} b_i \otimes c_i
\]

with at least two non-collinear pairs \((b_i, c_i)\). This is the same
mathematical distinction used in quantum cognition models of human
decision-making, where interference terms arise precisely when cognitive
state cannot be decomposed into independent ``category'' and
``decision'' components.

Patent-5's ``bivector interlingua'' can be understood as a particular
mapping out of this tensor space. Let \(V\) be a real vector space
equipped with a wedge product \(\wedge\), and let \(\wedge^2 V\) denote
its grade-2 subspace. A bivector interlingua is then a map

\[
\mathcal{B} : \mathcal{H}_{\text{behavior}} \otimes \mathcal{H}_{\text{context}} \to \wedge^2 V
\]

with two key properties:

\begin{enumerate}
\def\labelenumi{\arabic{enumi}.}
\tightlist
\item
  \textbf{Orientation preservation.} Projections that correspond to
  ``axes'' in the Plot of Time frame --- such as intelligence-community
  funding versus civilian surface, or artist-first versus state-first
  dynamics --- map to oriented basis pairs in \(V\).
\item
  \textbf{Script control.} User-level scripts correspond to choices of
  projection from \(\wedge^2 V\) into different rendering bases
  (biological, financial, narrative) without altering the underlying
  bivector.
\end{enumerate}

Intuitively, \(\mathcal{B}\) takes a combined behavioral--context state
\(\Psi\) and encodes it as an oriented ``area'' in an abstract space.
Different observables --- benchmark scores, credit flows, historical
arcs --- are different slices through that same bivector. The
non-separable structure is what allows interference: when two persona
contexts are applied in sequence, their associated projections need not
commute, and the resulting composite can perform worse than either
alone, as the quantum-cognition analysis of multi-persona prompts
suggests.

\begin{quote}
\textbf{Claim 7.2 (Non-separability requirement for interference).} Any
routing architecture that treats behavioral state and context as
strictly separable --- i.e., that always factorizes session state into
\(b \otimes c\) and bases all decisions on \(b\) alone --- cannot
represent the destructive interference effects observed when multiple
persona contexts are stacked. In such an architecture, additional
context can at worst be ignored; it cannot systematically make a
high-fidelity context worse.
\end{quote}

The substrate does not attempt to compute interference terms explicitly
at this stage. It does, however, choose a representational form ---
tensor product followed by an antisymmetric map into \(\wedge^2 V\) ---
that is capable of hosting them once the routing layer is upgraded.

\subsection{PACT, invariants, and system-level
implications}\label{pact-invariants-and-system-level-implications}

Putting these pieces together, a PACT-driven session can be viewed as a
sequence of linear operators \(\{U_t\}\) acting on combined state:

\[
\Psi_{t+1} = U_t(\Psi_t),
\]

where each \(U_t\) is induced by a specific PACT envelope and its
resulting closure: it updates the request header in \(X\), performs a
model call subject to budget and tier constraints, and maps the new
behavioral--context pair into a closure \(C_t\) via the BCODE operator
and, conceptually, into a bivector via \(\mathcal{B}\).

Three invariants then characterize the substrate's behavior in this
simplified view:

\begin{enumerate}
\def\labelenumi{\arabic{enumi}.}
\tightlist
\item
  \textbf{Budget monotonicity.} Credit balances recorded in PACT headers
  are non-increasing over time:
\end{enumerate}

\[
B_{t+1} \le B_t,
\]

with equality holding only for steps that do not trigger model compute.

\begin{enumerate}
\def\labelenumi{\arabic{enumi}.}
\setcounter{enumi}{1}
\item
  \textbf{Audit monotonicity.} The multiset of closures \(\{C_i\}\) is
  non-decreasing: each PACT application appends exactly one closure to
  the log, accepted or rejected, and no existing closure is modified or
  removed.
\item
  \textbf{Closure band constraint.} Accepted closures satisfy
  \(\|C_i\|_F \le d\) for a configured threshold \(d\), and any closure
  with \(\|C_i\|_F > d\) is explicitly marked as rejected.
\end{enumerate}

In combination, these invariants give the substrate a property that is
unusual in current AI systems: \textbf{behavioral state is both
non-separable and audit-constrained.} The system's state includes
entangled behavior--context bivectors that are capable in principle of
supporting interference and phase transitions, but every step in its
evolution is recorded as a signed closure with bounded norm and monotone
credit consumption.

From a research perspective, this suggests two follow-on directions.

First, a formal non-separability index (NSI) should be defined at the
level of \(\Psi_t\) or its image under \(\mathcal{B}\), not at the level
of single-model outputs. The ``stripped claims'' around NSI in earlier
drafts reflect the absence of such a definition; the present
specification creates the mathematical space in which it could be
meaningfully introduced but does not yet propose a concrete estimator.

Second, routing policies can be treated as explicit operators on
\(\Psi_t\) rather than as ad-hoc heuristic rules. A routing policy that
respects non-separability would choose actions based not only on
marginal persona fidelity or cost, but on how different PACT
applications are expected to move \(\Psi_t\) within
\(\mathcal{H}_{\text{behavior}} \otimes \mathcal{H}_{\text{context}}\),
including the possibility of interference when incompatible projections
are composed.

The substrate therefore occupies an unusual point in the design space.
It is not yet a full topological-intelligence system, but it already
treats session state as a non-separable object evolved by typed
operators, under explicit budget and audit invariants. This is
sufficient to support the empirical work in Sections 3 and 4, the safety
guarantees in Section 5, and the tier ladder in Section 6, while leaving
room for future metrics and routing behaviors that quantify and exploit
non-separability directly.

\section{Open artifacts and
replication}\label{open-artifacts-and-replication}

To support verification and reuse, the work is published with a concrete
artifact set rather than only narrative claims. The paper itself is
produced from Markdown (\texttt{beyond-web5-substrate-map.md}) using a
standard toolchain (Pandoc + the NeurIPS 2026 LaTeX template) so that
section numbering, citations, and tables remain stable across formats.
The intent is that a technically inclined reader can, with modest
effort, (i) audit or rerun the ConstellationBench benchmark via the
replication recipe in §8.2, (ii) stand up the same minimal state service
and wire tests via the execution receipt in §4, and (iii) inspect or
adapt the protocol and pricing scaffolds for their own multi-agent
systems.

\subsection{Open artifacts}\label{open-artifacts}

What is publicly available now: the ConstellationBench dataset on
Hugging Face (results JSON + signed bivector sidecars +
\texttt{eval.yaml}); the BCODE module (closure operator, replay
verifier, unit tests); the protocol specification accompanying this
paper; and the Phase-0 state service with its wire-test harness.
Anything that touches measurements (dataset, code, configuration) is
released under MIT; business-model documents (Constellation Credits
design, \$MOTION economics) are held more tightly while the bridge to
public on-chain mechanics is finalized.

{\def\LTcaptype{none} % do not increment counter
\begin{longtable}[]{@{}
  >{\raggedright\arraybackslash}p{(\linewidth - 8\tabcolsep) * \real{0.2000}}
  >{\raggedright\arraybackslash}p{(\linewidth - 8\tabcolsep) * \real{0.2000}}
  >{\raggedright\arraybackslash}p{(\linewidth - 8\tabcolsep) * \real{0.2000}}
  >{\raggedright\arraybackslash}p{(\linewidth - 8\tabcolsep) * \real{0.2000}}
  >{\raggedright\arraybackslash}p{(\linewidth - 8\tabcolsep) * \real{0.2000}}@{}}
\toprule\noalign{}
\begin{minipage}[b]{\linewidth}\raggedright
Artifact
\end{minipage} & \begin{minipage}[b]{\linewidth}\raggedright
Location
\end{minipage} & \begin{minipage}[b]{\linewidth}\raggedright
Contents
\end{minipage} & \begin{minipage}[b]{\linewidth}\raggedright
License
\end{minipage} & \begin{minipage}[b]{\linewidth}\raggedright
Cost to reuse
\end{minipage} \\
\midrule\noalign{}
\endhead
\bottomrule\noalign{}
\endlastfoot
ConstellationBench dataset &
huggingface.co/datasets/AirlockLabs/constellation-bench & JSONL results,
signed bivector sidecars, \texttt{eval.yaml}, README & MIT & \$0 (cached
transcripts; no API calls needed) \\
ConstellationBench leaderboard &
huggingface.co/spaces/AirlockLabs/constellation-bench-leaderboard &
Interactive scoring + replay verifier & MIT & \$0 \\
BCODE module & \texttt{constellation-bench-hf/bcode/} &
\texttt{closure.py}, \texttt{verify.py}, \texttt{test\_closure.py},
\texttt{fidelity.py} & MIT & \$0 (local CPU only) \\
Protocol specification & This document + \texttt{protocol-spec.md} &
Roster, tasks, config, run list, stripped-claim list & MIT & \$0 \\
State service (Phase-0) & \texttt{airlock-coordination/server.py} +
\texttt{validate-wires.py} & FastAPI + SQLAlchemy + Postgres; pinned
\texttt{requirements.txt} & MIT & \$0 (Python venv; \textasciitilde5 min
setup) \\
Pricing + credits design &
\texttt{2026-03-11-constellation-credits-*.md} & CC unit spec, ledger
scheme, dial behavior & Internal (held during economics finalization) &
N/A \\
Spacebar mobile client & (forthcoming --- humbble fork) & RN+Expo native
shell + Skeet bot federation surface & MIT & \$0 (Apple/Google dev
account costs at deploy time) \\
\end{longtable}
}

These artifacts pin model slugs, temperature, max tokens, and task
counts, so future researchers can detect drift even when exact replicas
become expensive (for example, if a model is deprecated or a provider
changes pricing). The pinning is what makes the dataset useful as a
longitudinal benchmark, not just a snapshot.

\subsection{Benchmark replication
recipe}\label{benchmark-replication-recipe}

This subsection covers benchmark replication --- re-running or auditing
the ConstellationBench dataset --- which is conceptually distinct from
the system execution receipt in §4 (bringing up the state service). A
full ConstellationBench rerun is approximately 22,200 calls /
\textasciitilde\$115 USD on OpenRouter at the prices available at time
of writing; a smoke test (1 persona, 1 model, 7 benchmarks) is
approximately 50 calls / under \$1. The recommended path is to read the
cached dataset first, run the verifier against it, and only execute new
API calls if independent confirmation is needed.

A minimal reproduction looks like:

\begin{enumerate}
\def\labelenumi{\arabic{enumi}.}
\tightlist
\item
  \textbf{Clone / download dataset and repo.}

  \begin{itemize}
  \tightlist
  \item
    \texttt{huggingface-cli\ download\ AirlockLabs/constellation-bench\ -\/-repo-type\ dataset}
  \item
    \texttt{git\ clone} the repository containing \texttt{bcode/} and
    the \texttt{scripts/} directory.
  \end{itemize}
\item
  \textbf{Install dependencies.}

  \begin{itemize}
  \tightlist
  \item
    \texttt{pip\ install\ -r\ requirements.txt} (pinned versions of
    FastAPI, SQLAlchemy, pydantic, ULID, PyYAML, plus the BCODE/DECF
    libraries).
  \end{itemize}
\item
  \textbf{Smoke test run.}

  \begin{itemize}
  \tightlist
  \item
    \texttt{python\ scripts/quick\_bench.py\ -\/-model\ haiku-4.5\ -\/-persona\ maverick}
  \item
    Expected: \textasciitilde7 benchmarks, \textasciitilde50 calls,
    completes in a few minutes; produces a sidecar JSONL file with
    signed bivector closures.
  \end{itemize}
\item
  \textbf{Full protocol run} (optional).

  \begin{itemize}
  \tightlist
  \item
    \texttt{python\ scripts/quick\_bench.py\ -\/-full}
  \item
    Expected: 22 models, 17 personas, 7 benchmarks,
    \textasciitilde22,200 calls, \textasciitilde\$115 USD on OpenRouter.
  \end{itemize}
\item
  \textbf{Verify sidecars.}

  \begin{itemize}
  \tightlist
  \item
    \texttt{python\ -m\ bcode.verify\ results-\textless{}model-slug\textgreater{}-bivectors.jsonl}
  \item
    Expected output: \texttt{VERIFIED\ N,\ MISMATCH\ 0,\ BAD-SIG\ 0}
    with N equal to the line count in the sidecar file. Any non-zero
    MISMATCH or BAD-SIG indicates either tampering or replay drift.
  \end{itemize}
\item
  \textbf{Recompute summary statistics.}

  \begin{itemize}
  \tightlist
  \item
    From the JSONL results, recompute composite scores and leaderboards
    using the same scoring rubric included in the dataset README. If
    aggregate numbers diverge from those reported here by more than the
    documented variance, suspect model-version drift at the provider
    level.
  \end{itemize}
\end{enumerate}

A note on limits: model APIs are stateless and frontier providers do not
guarantee reproducibility across versions, so perfect deterministic
replay is impossible. Reproducibility is achieved via cached transcripts
and sidecar verification, not via seed-level determinism. The intent is
that any reader can verify the BCODE closures, recompute the scoring,
and check for drift; running new calls is optional, not required, to
confirm the paper's claims.

\subsection{Derivation map: this paper as one slice of a master
spine}\label{derivation-map-this-paper-as-one-slice-of-a-master-spine}

This paper is one derivation of a denser internal master document. The
same spine is intended to feed product requirement documents, investor
decks, grant applications, editorial essays, and onboarding runbooks;
the paper's role is to be the academically defensible slice that other
audiences can extend or compress without re-deriving canon. Readers who
wish to use this work for purposes beyond reading can therefore draw
from the following derivation map:

\begin{itemize}
\tightlist
\item
  \textbf{Engineers building on the substrate.} Take §§4, 6, and 7.1 as
  the operational specification for an open-source ConstellationBench
  implementation, BCODE/PACT service, and tier ladder. The state-service
  architecture in §4 maps directly to a service-boundary PRD; the typed
  PACT operator in §7.1 specifies the protocol envelope schema; §6's
  tier-ladder bundling specifies pricing-page and credit-ledger
  behavior. The narrative and speculative-theory portions are not
  load-bearing for engineering work.
\item
  \textbf{Operators and product teams.} Take §2 (substrate map), §6
  (tier ladder), and §8 (artifact map) as the starting point for PRDs
  and service boundaries. The paper is system specification, not only
  benchmark report; the seven primitives in Table 1 each correspond to a
  buildable component with a documented role.
\item
  \textbf{Researchers extending the empirical work.} Take §§3-5 and §7.2
  as the empirical and theoretical wedges. Use the dataset and code to
  run new routing experiments, interference tests, or non-separability
  metrics. The ``stripped claims'' inventory in the protocol
  specification identifies exactly which earlier research-grade claims
  have been removed pending evidence; those represent open research
  directions.
\item
  \textbf{Funders, grant reviewers, and strategic partners.} Take the
  abstract + §§1-2 + §6 + §8 as the project overview and reproducibility
  posture. The same spine can be attached to grant applications,
  investor decks, and editorial essays without rewriting facts;
  cross-audience consistency is a feature.
\end{itemize}

The intent is that a single, evidence-aligned spine supports multiple
public artifacts --- this paper, the dataset, downstream PRDs, and decks
--- so that no audience is asked to trust a claim that is not grounded
in at least one open artifact.

\subsection{References}\label{references}

\begin{itemize}
\tightlist
\item
  Berners-Lee, T. (1989). \emph{Information Management: A Proposal.}
  CERN. \texttt{{[}@bernerslee1989{]}}
\item
  Berners-Lee, T. et al.~(2016). \emph{Solid: A Decentralized Protocol
  for Personal Data Stores on the Web.} MIT CSAIL.
  \texttt{{[}@solid\_pods{]}}
\item
  Dorsey, J., Block Inc.~(2022). \emph{TBD's Web5: A Decentralized Web
  Platform with Self-Sovereign Identity.} Block Inc.~Technical
  Specification. \texttt{{[}@tbd\_web5{]}}
\item
  Holwerda, Z. (2026). \emph{ConstellationBench: A 17-Persona Benchmark
  for Behavioral Fidelity in Frontier Language Models.} Airlock Labs
  Technical Report.
  https://huggingface.co/datasets/AirlockLabs/constellation-bench
  \texttt{{[}@constellationbench2026{]}}
\item
  Holwerda, Z. (2026). \emph{PACT Protocol Specification: A Typed
  Operator for Operator-Resident AI Substrates.} Airlock Labs Technical
  Report. \texttt{{[}@pact\_protocol\_spec{]}}
\item
  Holwerda, Z. (2026). \emph{Beyond Web5: Topological Intelligence and
  the Operator-Resident Substrate (white paper v7).} Airlock Labs Public
  White Paper. \texttt{{[}@white\_paper\_v7\_TI{]}}
\item
  Holwerda, Z. (2026). \emph{WP2: Entanglement-Safe Handshake --- A
  Behavioral Commitment Primitive.} Airlock Labs Whitepaper Series.
  \texttt{{[}@wp2\_entanglement\_safe\_handshake{]}}
\item
  Holwerda, Z. (2026). \emph{WP3: SHA-256 as a Layered Trust Object ---
  A Five-Layer Frame for Behavioral Commitments.} Airlock Labs
  Whitepaper Series. \texttt{{[}@wp3\_sha256\_layered\_trust{]}}
\item
  Sporns, O., O'Connor, J., et al.~(2023). \emph{Topological
  Intelligence: From Substrate-Resident Agents to Operator-Resident
  Systems.} (Survey of Web4-tier substrate literature.)
  \texttt{{[}@oconnor2023web4{]}}
\item
  Katsarosgiannis, G. (2026, April 29). \emph{Stern-Brocot
  Reconstruction of Anthropic's Subscription Credit Caps.} Public
  disclosure (Threads post). https://threads.net/@katsarosgiannis
  \texttt{{[}@katsarosgiannis2026{]}}
\item
  HiddenLayer Inc.~(2024). \emph{MLSec Platform: Adversarial-Attack
  Detection and Behavioral Anomaly Monitoring for Production Models.}
  Commercial product documentation. https://hiddenlayer.com
  \texttt{{[}@hiddenlayer{]}}
\item
  Anthropic. (2024). \emph{Model Context Protocol (MCP) Specification.}
  https://modelcontextprotocol.io \texttt{{[}@mcp\_anthropic{]}}
\item
  Zhang, A., Lipton, Z., Li, M., Smola, A. (2024). \emph{Dive into Deep
  Learning.} d2l.ai. https://d2l.ai \texttt{{[}@d2l\_book{]}}
\item
  Mühle, A., Grüner, A., Gayvoronskaya, T., Meinel, C. (2018). \emph{A
  Survey on Essential Components of a Self-Sovereign Identity.} Computer
  Science Review, 30, 80-86. \texttt{{[}@ssi\_survey{]}}
\item
  W3C. (2022). \emph{Decentralized Identifiers (DIDs) v1.0: Core
  Architecture, Data Model, and Representations.} W3C Recommendation.
  https://www.w3.org/TR/did-core/ \texttt{{[}@w3c\_did\_core{]}}
\end{itemize}

\begin{center}\rule{0.5\linewidth}{0.5pt}\end{center}

\textbf{Draft notes for editor:}

\begin{itemize}
\tightlist
\item
  Abstract + §1 + §2 + Table 1 are operator-locked (see canon
  \texttt{project\_neurips\_paper\_pivot\_substrate\_map\_beyond\_web5.md});
  use verbatim.
\item
  §3 prose to be migrated from \texttt{white-paper-v3.md} (operator's
  narrative draft) with translation to academic register per
  \texttt{reference\_white\_paper\_v3\_love\_fourth\_dimension\_constellationbench\_pact.md}.
\item
  §4 prose to follow once Phase-0 wiring details are written up (state
  service architecture, validate-wires.py output, bootstrap procedure).
\item
  §5 prose to draw from
  \texttt{reference\_stern\_brocot\_side\_channel\_pricing\_recovery.md}
  + \texttt{reference\_hiddenlayer\_gatekept\_web5\_security\_slot.md}.
\item
  §6 prose to draw from
  \texttt{project\_skeet\_mail\_phone\_anchored\_burner\_identity.md}
  D2L ladder section.
\item
  §7 prose to draw from
  \texttt{feedback\_distribute\_dont\_consolidate\_ip\_cao\_posture.md},
  \texttt{feedback\_warp\_drive\_not\_wrap\_feathering\_ip.md},
  white-paper-v3.md PACT section,
  \texttt{feedback\_chainlink\_skeet\_integration\_phased.md}.
\item
  §8 reproducibility table is current as of 2026-04-29; verify all links
  and costs before submission.
\item
  Citation style: Pandoc-flavored Markdown (\texttt{{[}@key{]}}); LaTeX
  export will map to \texttt{\textbackslash{}cite\{key\}}.
\item
  LaTeX template at \texttt{papers/latex-neurips/neurips\_2026.tex} and
  \texttt{neurips\_2026.sty} is already in place; \texttt{paper.tex}
  will receive the export.
\item
  Page-limit risk: §3 + §7 are now significantly larger than originally
  outlined (white-paper-v3.md migration adds substantial content). Tight
  editing pass on May 3 critical.
\end{itemize}
